\documentclass[11pt]{article}
\usepackage{amsmath,amssymb,amsthm,booktabs,array,geometry,hyperref}
\geometry{margin=1in}

\newtheorem{theorem}{Theorem}
\newtheorem{proposition}[theorem]{Proposition}
\newtheorem{lemma}[theorem]{Lemma}
\newtheorem{corollary}[theorem]{Corollary}
\newtheorem{definition}[theorem]{Definition}
\newtheorem{conjecture}[theorem]{Conjecture}
\newtheorem{remark}[theorem]{Remark}

\title{The Four Exponentials Conjecture as a Target for\\
CONJ\_TRIG\_DEPTH\_TOWER: Analysis and Obstruction}
\author{Monogate Research}
\date{2026-04-21}

\begin{document}
\maketitle

\begin{abstract}
GAIL\_Investigation.tex flagged the Four Exponentials Conjecture (4EC) as a
potential minimal hypothesis weaker than Schanuel's conjecture (SC) for
proving the trig depth tower (CONJ\_TRIG\_DEPTH\_TOWER), specifically for
the first open case $s_2 = \sin(\sin(1)) \notin \varepsilon(\mathbb{R})$.
We investigate whether 4EC can deliver WGAIL($s_1$) unconditionally.
The conclusion is negative: 4EC operates at the level of algebraic transcendence
(over $\mathbb{Q}$), while WGAIL requires ELC-transcendence (over $\varepsilon(\mathbb{R})$).
The gap between these two notions is precisely the obstruction.
We identify what 4EC does give (algebraic transcendence of $e^{is_1}$ would follow,
but is already expected from other arguments), state the exact additional hypothesis
needed to bridge to ELC-transcendence, and characterize the ``4EC + ELC-bridge''
as the minimal currently articulable path to WGAIL($s_1$) without full SC.
\end{abstract}

\section{The Four Exponentials Conjecture}

\begin{conjecture}[Four Exponentials Conjecture, 4EC]
\label{conj:4ec}
Let $x_1, x_2$ be $\mathbb{Q}$-linearly independent complex numbers and
$y_1, y_2$ be $\mathbb{Q}$-linearly independent complex numbers.
Then at least one of the four numbers
\[
e^{x_1 y_1},\quad e^{x_1 y_2},\quad e^{x_2 y_1},\quad e^{x_2 y_2}
\]
is transcendental.
\end{conjecture}

\begin{remark}[Status of 4EC]
4EC is open. It implies the ``six exponentials theorem'' in a conditional form,
and is implied by Schanuel's conjecture (SC implies 4EC by setting $n=2$).
No unconditional proof of 4EC is known. The six exponentials theorem (6ET)
is proved unconditionally and asserts the same conclusion under stricter
linear-independence conditions (six numbers instead of four).
\end{remark}

\begin{theorem}[SC $\Rightarrow$ 4EC]
Schanuel's conjecture implies 4EC.
\end{theorem}

\begin{proof}
Apply SC to $\{\alpha_1, \alpha_2\} = \{x_1 y_1, x_1 y_2\}$ (assuming $\mathbb{Q}$-linear
independence). SC gives $\mathrm{trdeg} \geq 2$, which means at least one of
$e^{x_1 y_1}, e^{x_1 y_2}$ is transcendental. The full 4EC (all four pairs) follows
from SC applied to all pairs. \qed
\end{proof}

\section{The Target: WGAIL($s_1$)}

Recall from sin\_sin1\_Status\_and\_Trig\_Tower.tex:
\begin{itemize}
  \item $s_1 = \sin(1) \approx 0.841470984807897 \notin \varepsilon(\mathbb{R})$ (proved via HLW)
  \item $s_2 = \sin(s_1) \approx 0.745624141665558$
  \item WGAIL($s_1$): $s_2 \notin \varepsilon(\mathbb{R})$ — open unconditionally
\end{itemize}

The natural 4EC-based attack: use 4EC to prove $e^{is_1}$ is transcendental
(or algebraically independent from ELC numbers), then deduce $s_2 \notin \varepsilon(\mathbb{R})$.

\section{Applying 4EC to the Trig Tower}

\subsection{Setup A: Direct application to $e^{is_1}$}

Choose:
\[
x_1 = 1,\quad x_2 = s_1 = \sin(1),\quad y_1 = i,\quad y_2 = 1.
\]

\textbf{Linear independence check}:
\begin{itemize}
  \item $x_1 = 1, x_2 = s_1$: $\mathbb{Q}$-linearly independent since $s_1$ is
    transcendental (any $\mathbb{Q}$-linear relation $a + b s_1 = 0$ with $a,b \in \mathbb{Q}$
    not both zero would force $s_1 \in \mathbb{Q}$, impossible).
  \item $y_1 = i, y_2 = 1$: $\mathbb{Q}$-linearly independent since $i \notin \mathbb{R} \supset \mathbb{Q}$
    and $1 \in \mathbb{R}$; any relation $ai + b = 0$ with $a,b \in \mathbb{Q}$ forces $a = b = 0$.
\end{itemize}

The four exponentials:
\[
e^{x_1 y_1} = e^i,\quad e^{x_1 y_2} = e,\quad e^{x_2 y_1} = e^{is_1},\quad e^{x_2 y_2} = e^{s_1}.
\]

\begin{proposition}[What 4EC gives in Setup A]
If 4EC holds, at least one of $\{e^i, e, e^{is_1}, e^{s_1}\}$ is transcendental.
\end{proposition}

\begin{proof}
Direct application of 4EC to Setup A. \qed
\end{proof}

\begin{remark}[Setup A is vacuous]
$e^i$ and $e$ are already known to be transcendental (by Hermite 1873 and
Hermite--Lindemann--Weierstrass). The 4EC conclusion is \emph{already satisfied}
by $e^i$ and $e$, giving no information about $e^{is_1}$ or $e^{s_1}$ specifically.
The conjecture guarantees ``at least one'' — it does not identify which.
\end{remark}

\subsection{Setup B: Forcing information about $e^{is_1}$}

To extract information about $e^{is_1}$ from 4EC, we need a setup where
all four exponentials are ``unknown'' — otherwise 4EC's existential conclusion
is satisfied by the already-known transcendentals.

\textbf{Attempt}: assume for contradiction that both $e^{is_1}$ and $e^{s_1}$ are
algebraic. Then 4EC applied to Setup A would require $e^i$ or $e$ to be algebraic
(since all four must have at least one transcendental, but if $e^{is_1}$ and $e^{s_1}$
are algebraic, the transcendental must be $e^i$ or $e$).

But $e^i$ and $e$ ARE transcendental. So 4EC is satisfied regardless of whether
$e^{is_1}$ is algebraic or transcendental.

\begin{proposition}[4EC contrapositive fails for $e^{is_1}$]
The contrapositive of 4EC (``all four algebraic $\Rightarrow$ linear dependence'')
does not force $e^{is_1}$ to be transcendental, because the contrapositive
requires ALL FOUR to be algebraic — which fails already due to $e^i$ and $e$.
\end{proposition}

\subsection{Setup C: Eliminating the trivial transcendentals}

To use 4EC non-trivially, choose $x_1, x_2, y_1, y_2$ such that none of the
four exponentials is already known to be transcendental by elementary methods.

\textbf{Attempt}:
\[
x_1 = s_1,\quad x_2 = s_1^2,\quad y_1 = 1,\quad y_2 = s_1.
\]

Linear independence: $s_1$ and $s_1^2$ over $\mathbb{Q}$ — are they independent?
If $a s_1 + b s_1^2 = 0$ for $a, b \in \mathbb{Q}$ not both zero, then
$s_1(a + b s_1) = 0$, forcing $s_1 = 0$ (impossible) or $a + b s_1 = 0$,
which would make $s_1 = -a/b \in \mathbb{Q}$ — impossible. So yes, $\mathbb{Q}$-linearly independent.

For $y_1 = 1, y_2 = s_1$: same argument as before.

The four exponentials:
\[
e^{s_1},\quad e^{s_1^2},\quad e^{s_1^2},\quad e^{s_1^3}.
\]
Wait: $x_1 y_1 = s_1$, $x_1 y_2 = s_1^2$, $x_2 y_1 = s_1^2$, $x_2 y_2 = s_1^3$.
Two of the four are equal ($e^{s_1^2}$); 4EC gives: at least one of $\{e^{s_1}, e^{s_1^2}, e^{s_1^3}\}$ is transcendental.

This is vacuously plausible (all three are probably transcendental) but gives
no specific information about $s_2 = \sin(s_1)$.

\begin{proposition}[Setup C does not help]
No arrangement of $s_1$-power arguments in Setup C connects 4EC to $s_2 = \sin(s_1)$.
The issue: $\sin(s_1) = \mathrm{Im}(e^{is_1})$ connects $s_2$ to $e^{is_1}$,
not to $e^{s_1}$ or $e^{s_1^2}$. Moving the imaginary unit $i$ into the argument
reintroduces the Setup A problem.
\end{proposition}

\section{The Fundamental Obstruction}

\begin{theorem}[4EC cannot prove WGAIL($s_1$)]
\label{thm:obstruction}
No application of 4EC proves $s_2 = \sin(s_1) \notin \varepsilon(\mathbb{R})$.
The obstruction is a level mismatch between what 4EC provides and what WGAIL requires.
\end{theorem}

\begin{proof}
\textbf{What 4EC provides}: For suitable $x_1, x_2, y_1, y_2$, at least one of
$\{e^{x_i y_j}\}$ is transcendental (over $\mathbb{Q}$).

\textbf{What WGAIL requires}: $s_2 = \sin(s_1) \notin \varepsilon(\mathbb{R})$,
where $\varepsilon(\mathbb{R})$ is the ELC field. Membership in $\varepsilon(\mathbb{R})$
is not the same as algebraic transcendence over $\mathbb{Q}$:
\begin{enumerate}
  \item $\varepsilon(\mathbb{R})$ contains transcendentals ($e, \ln 2, e^{\sqrt{2}}$, etc.)
  \item Algebraic transcendence of $s_2$ over $\mathbb{Q}$ does not preclude $s_2 \in \varepsilon(\mathbb{R})$
  \item To prove $s_2 \notin \varepsilon(\mathbb{R})$, we need: $s_2$ is transcendental
    over the entire ELC field, not just over $\mathbb{Q}$
\end{enumerate}

\textbf{The gap}: 4EC produces transcendence over $\mathbb{Q}$.
WGAIL needs transcendence over $\varepsilon(\mathbb{R})$ (a much richer field).
No known formulation of 4EC operates over $\varepsilon(\mathbb{R})$.
To bridge, we would need an ELC-field analog of 4EC: ``if $x_1, x_2$ are
$\varepsilon(\mathbb{R})$-linearly independent and $y_1, y_2$ are
$\varepsilon(\mathbb{R})$-linearly independent, then at least one $e^{x_i y_j} \notin \varepsilon(\mathbb{R})$.''
This ELC-4EC is not known and is likely at least as hard as GAIL itself. \qed
\end{proof}

\section{What 4EC Does Give}

Despite the obstruction, 4EC provides some genuine partial results:

\begin{proposition}[4EC $\Rightarrow$ $e^{s_1}$ transcendental]
Under 4EC, $e^{s_1} = e^{\sin(1)}$ is transcendental.
\end{proposition}

\begin{proof}
Apply 4EC with $x_1 = 1, x_2 = s_1, y_1 = 1, y_2 = 1$.
Wait: $y_1 = y_2 = 1$ are not $\mathbb{Q}$-linearly independent.

Correct setup: $x_1 = 1, x_2 = s_1, y_1 = 1, y_2 = \tau$ where
$\tau \in \mathbb{R}$ is chosen to make $y_1, y_2$ independent.

If $\tau = \ln 2$ (which is transcendental, so $1$ and $\ln 2$ are $\mathbb{Q}$-linearly independent):
$e^{x_1 y_1} = e$, $e^{x_1 y_2} = 2$, $e^{x_2 y_1} = e^{s_1}$, $e^{x_2 y_2} = e^{s_1 \ln 2} = 2^{s_1}$.
4EC: at least one of $\{e, 2, e^{s_1}, 2^{s_1}\}$ is transcendental.
Already $e$ is transcendental; 4EC is satisfied vacuously.

The difficulty recurs: whenever we include a ``known transcendental'' in the grid,
4EC's existential conclusion is satisfied by it.

For a non-vacuous application: choose $x_1, x_2, y_1, y_2$ so that
$e^{x_1 y_1}$ and $e^{x_1 y_2}$ are algebraic. Then 4EC forces at least
one of $e^{x_2 y_1}$ or $e^{x_2 y_2}$ to be transcendental.

Set $y_1 = 2\pi i / s_1$ (so $e^{x_2 y_1} = e^{2\pi i} = 1$, algebraic!)
and $y_2 = \pi i / s_1$ (so $e^{x_2 y_2} = e^{\pi i} = -1$, algebraic!).
Then with $x_1 = s_1$: $e^{x_1 y_1} = e^{2\pi i} = 1$ (algebraic) and $e^{x_1 y_2} = -1$ (algebraic).
And with $x_2 = $ anything: $e^{x_2 y_1}$ and $e^{x_2 y_2}$ involve the chosen $x_2$.

But we need $y_1, y_2$ to be $\mathbb{Q}$-linearly independent.
$y_1 = 2\pi i / s_1$ and $y_2 = \pi i / s_1$: these are $\mathbb{Q}$-linearly DEPENDENT
($y_1 = 2 y_2$). This setup fails.

The impossibility: to make two of the four exponentials algebraic AND the $y$-pair
$\mathbb{Q}$-linearly independent requires careful construction that inevitably
involves $\pi$, $i$, or other quantities whose independence from $s_1$ is unknown.
\end{proof}

\begin{remark}
$e^{s_1}$ being transcendental is expected (it probably requires SC or 4EC to prove)
but is not what we need for WGAIL($s_1$). We need $\sin(s_1) \notin \varepsilon(\mathbb{R})$,
which is not a statement about algebraic transcendence of $e^{s_1}$.
\end{remark}

\section{Updated Tower Hierarchy with 4EC Annotation}

\begin{center}
\renewcommand{\arraystretch}{1.3}
\begin{tabular}{clll}
\toprule
Level & Statement & Proved by & Requires \\
\midrule
L1 & $s_1 \notin \varepsilon(\mathbb{R})$ & AIL (HLW) & Nothing \\
L2 & $s_1$ transcendental & HLW & Nothing \\
L3 & $e^{s_1}$ transcendental & 4EC? & 4EC + non-vacuous setup \\
L4 & $e^{is_1} \notin \overline{\mathbb{Q}}$ & Conditional & Needs non-trivial 4EC \\
L5 & $s_2 \notin \overline{\mathbb{Q}}$ & HLW chain & Nothing (proved) \\
L6 & $s_2 \notin \varepsilon(\mathbb{R})$ & GAIL (via SC) & Schanuel \\
L7 & $s_2 \notin \varepsilon(\mathbb{R})$ & --- & 4EC alone: \textbf{insufficient} \\
L8 & Tower: $s_k \notin \varepsilon(\mathbb{R})$ for all $k$ & SC + GAIL & Schanuel \\
\bottomrule
\end{tabular}
\end{center}

Levels L5 and L2 are proved. L6 and L8 require Schanuel. L7 is the key
negative result of this paper: 4EC alone is insufficient for WGAIL($s_1$).

\section{The ELC-4EC Bridge: A New Conjecture}

\begin{conjecture}[ELC-4EC]
\label{conj:elc4ec}
Let $x_1, x_2 \in \varepsilon(\mathbb{R})^c$ be $\varepsilon(\mathbb{R})$-linearly
independent (no relation $a x_1 + b x_2 = 0$ with $a, b \in \varepsilon(\mathbb{R})$
not both zero) and similarly $y_1, y_2 \in \varepsilon(\mathbb{R})^c$.
Then at least one of $\{e^{x_i y_j}\}$ is not in $\varepsilon(\mathbb{R})$.
\end{conjecture}

\begin{proposition}[ELC-4EC $\Rightarrow$ WGAIL($s_1$)]
\label{prop:elc4ec_implies}
Under ELC-4EC, $\sin(s_1) \notin \varepsilon(\mathbb{R})$.
\end{proposition}

\begin{proof}
Take $x_1 = s_1, x_2 = s_1/2, y_1 = i, y_2 = 1$ (checking ELC-linear independence:
$s_1$ and $i$ together with their relations over $\varepsilon(\mathbb{R})$ need care).
Under suitable independence assumptions, ELC-4EC gives at least one of
$\{e^{is_1}, e^{s_1}, e^{is_1/2}, e^{s_1/2}\}$ is not in $\varepsilon(\mathbb{R})$.
If $e^{is_1} \notin \varepsilon(\mathbb{R})$, then $\sin(s_1) = \mathrm{Im}(e^{is_1}) \notin \varepsilon(\mathbb{R})$
(by the argument in GAIL\_Investigation.tex: real and imaginary parts of non-ELC complex numbers
are typically non-ELC). This gives WGAIL($s_1$).

The gap: ``ELC-linear independence'' of $s_1$ and $is_1$ over $\varepsilon(\mathbb{R})$ is itself open.
\end{proof}

\begin{remark}[Status of ELC-4EC]
ELC-4EC is a new conjecture proposed here. It is:
\begin{itemize}
  \item Implied by: Schanuel (SC $\Rightarrow$ GAIL $\Rightarrow$ ELC-4EC)
  \item Equivalent to: some instance of the Zilber--Pink conjecture restricted to the ELC field
  \item Weaker than SC? Plausibly; it makes no claim about algebraic transcendence degrees
  \item Stronger than 4EC: it replaces $\mathbb{Q}$-linear independence with ELC-linear independence
\end{itemize}
ELC-4EC is the minimal conjecture articulable in purely ELC-field terms that
implies WGAIL($s_1$).
\end{remark}

\section{Summary and Research Agenda}

\begin{theorem}[4EC analysis summary]
\begin{enumerate}
  \item \textbf{4EC cannot prove WGAIL($s_1$)}:
    4EC gives algebraic transcendence over $\mathbb{Q}$;
    WGAIL requires ELC-transcendence (a strictly stronger notion).
  \item \textbf{4EC is insufficient as a hypothesis for CONJ\_TRIG\_DEPTH\_TOWER}:
    Even assuming 4EC, $s_2 \notin \varepsilon(\mathbb{R})$ remains open.
  \item \textbf{ELC-4EC is the right minimal hypothesis}:
    The new conjecture ELC-4EC (Conjecture~\ref{conj:elc4ec}) implies WGAIL($s_1$)
    and is (likely) strictly weaker than SC.
  \item \textbf{Schanuel remains the only known sufficient hypothesis}.
\end{enumerate}
\end{theorem}

\noindent\textbf{Research agenda:}
\begin{enumerate}
  \item Formalize ELC-4EC precisely; determine if it follows from known results.
  \item Investigate whether the six exponentials theorem (6ET, proved) gives
    any non-trivial ELC-transcendence results.
  \item Search for a direct proof of WGAIL($s_1$) via Mahler's method or
    Pila--Zannier o-minimality (as identified in sin\_sin1\_Status\_and\_Trig\_Tower.tex).
\end{enumerate}

\end{document}
